% Preamble for "Gujarati, Step by Step, Root by Root" (HL00 book format).
% XeLaTeX only. Uses the vendored static Noto Sans Gujarati (see _fonts/),
% loaded by relative path so local and CI builds are identical. Rendered
% font-only via Script=Gujarati (no polyglossia Gujarati module needed).

\usepackage{fontspec}
\setmainfont{Latin Modern Roman}
\setsansfont{Latin Modern Sans}
\setmonofont{Latin Modern Mono}

\usepackage{newunicodechar}
\newunicodechar{ṁ}{\.{m}}
\newunicodechar{ṉ}{\b{n}}
\newunicodechar{ṟ}{\b{r}}
\newunicodechar{ḻ}{\b{l}}
\newunicodechar{ḷ}{\d{l}}
\newunicodechar{ũ}{\~{u}}
\newunicodechar{ĩ}{\~{\i}}
\newunicodechar{ã}{\~{a}}

\usepackage[table]{xcolor}
\usepackage{tcolorbox}
\tcbuselibrary{skins,breakable}
\usepackage{booktabs}
\usepackage{longtable}
\usepackage{tabularx}
\usepackage{array}
\usepackage[hidelinks]{hyperref}
\usepackage{titlesec}

\usepackage{polyglossia}
\setmainlanguage{english}

% Vendored Gujarati font (../../_fonts/ from <lang>/book/); \gu{...} = inline
% Gujarati snippet.
\newfontfamily\gujaratifont[
  Path=../../_fonts/,
  Script=Gujarati,
  BoldFont=NotoSansGujarati-Static.ttf,
  ItalicFont=NotoSansGujarati-Static.ttf,
  BoldItalicFont=NotoSansGujarati-Static.ttf
]{NotoSansGujarati-Static.ttf}
\newcommand{\gu}[1]{{\gujaratifont #1}}

\hypersetup{
  pdftitle={Gujarati, Step by Step, Root by Root},
  pdfauthor={Adhithya Rajasekaran},
  pdfsubject={Etymology-driven Gujarati curriculum},
  pdfkeywords={Gujarati, Gandhi, Indo-Aryan, language learning}
}
\pdfstringdefDisableCommands{%
  \def\gu#1{#1}% Keep Gujarati text while dropping font-only presentation.
  \def\gujaratifont{}%
}

% Lesson callouts are deliberately kept together when they fit. Let short pages
% end naturally instead of stretching their vertical glue around those boxes.
\raggedbottom

\newtcolorbox{cousinweb}{
  breakable, skin=enhanced, colback=yellow!8, colframe=yellow!45!black,
  boxrule=0.5pt, arc=1mm, left=6pt, right=6pt, top=4pt, bottom=4pt,
  fonttitle=\bfseries, title={The word, taken apart --- its roots and family}
}
\newtcolorbox{culture}{
  breakable, skin=enhanced, colback=red!5, colframe=red!40!black,
  boxrule=0.5pt, arc=1mm, left=6pt, right=6pt, top=4pt, bottom=4pt,
  fonttitle=\bfseries, title={Why it's said this way}
}
\newtcolorbox{grammarlens}[1][]{
  breakable, skin=enhanced, colback=blue!5, colframe=blue!35!black,
  boxrule=0.5pt, arc=1mm, left=6pt, right=6pt, top=4pt, bottom=4pt,
  fonttitle=\bfseries, title={Grammar Lens}, #1
}
\newtcolorbox{sounds}{
  breakable, skin=enhanced, colback=green!6, colframe=green!30!black,
  boxrule=0.5pt, arc=1mm, left=6pt, right=6pt, top=4pt, bottom=4pt,
  fonttitle=\bfseries, title={The letters in this word}
}

\titleformat{\chapter}[display]
  {\normalfont\Large\bfseries}{Chapter \thechapter}{10pt}{\huge}
